\section{Conclusion}
\label{sec:conclusion}

We have assembled and evaluated a workflow that separates colloidal molecular-dynamics
configurations into spatial domains, represents each with a fixed species-aware descriptor,
detects domains belonging to structural families absent from its training atlas, and
compares those domains against an independently constructed reference library that it may
decline to match.

Under strict leave-one-family-out evaluation, in which every statistical quantity is re-fitted
without the held-out family, the detector ranks unseen domains above all retained ones for
each of the three families tested. An external candidate rejected by the detector is
identified by the matcher against a reference absent from the training atlas. When the
correct reference is removed from the library, all twenty-three affected domains abstain
rather than accept the nearest survivor, which sits at a distance typical of an unrelated
prototype.

Two findings temper this. Admitting provisional machine-generated labels into training
produced no measurable improvement and degraded threshold stability by a factor of three.
And the acceptance threshold, rather than the representation, is the least stable component:
the predeclared calibration rule rejects roughly a quarter of in-distribution validation
domains, which alternative rules reduce substantially on the same data.

The extension of the atlas to newly identified families is supported by the architecture but
is not demonstrated here, and we do not claim it.
